\enquestion{4.1}{
    Suppose that $f(x) = e^{-2x}$ for $0 < x$. Determine the following:
    \begin{center}
        \begin{tabular}{*{3}{C{0.3\textwidth}}}
            (a) $P(1 < X)$ & (b) $P(1 < X < 2.5)$ & \multirow{2}{*}{(e) $P(3 \leq X)$}\\
            (c) $P(X = 3)$ & (d) $P(X < 4)$ & \\
            \multicolumn{2}{l}{(f) Determine $x$ such that $P(x < X) = 0.20$.} \\
            \multicolumn{2}{l}{(g) Determine $x$ such that $P(X \le x) = 0.10$.}
        \end{tabular}
    \end{center}
}

\enquestion{4.7}{
    Suppose that $f(x) = 2x^2$ for $-1 < x < 1$. Determine the following:

    \begin{center}
        \begin{tabular}{*{2}{m{0.45\textwidth}}}
            (a) $P(0 < X)$ & (b) $P(0.5 < X)$ \\
            (c) $P(-0.5 \le X \le 0.5)$ & (d) $P(X < -2)$ \\
            (e) $P(X < 0\text{ or } X > -0.5)$ & (f) $x$ such that $P(x < X) = 0.05$.
        \end{tabular}
    \end{center}
}

\enquestion{4.11}{
    The probability density function of the length of a metal rod is $f(x) = 2$ for $2.3 < x < 2.9$ meters.
    \begin{enumerate}[label=(\alph*)]
        \item If the specifications for this process are from $2.25$ to $2.75$ meters, what proportion of rods fail to meet the specifications?
        \item Assume that the probability density function is $f(x) = 2$ for an interval of length $0.5$ meters. Over what value should the density be centered to achieve the greatest proportion of rods within specifications?
    \end{enumerate}
}

\enquestion{4.13}{
    Suppose that the cumulative distribution function of the random variable $X$ is
    \[
        F(x) = \begin{cases}
            0 & x < 0 \\
            0.2x & 0 \le x < 5 \\
            1 & 5 \le x
        \end{cases}
    \]
    Determine the following:

    \begin{center}
        \begin{tabular}{*{2}{C{0.45\textwidth}}}
            (a) $P(X < 2.8)$ & (b) $P(X > 1.5)$ \\
            (c) $P(X < -2)$ & (d) $P(X > 6)$
        \end{tabular}
    \end{center}
}

\enquestion{4.15}{
    Determine the cumulative distribution function given the following probability density function:
    \[
        f(x) = e^{-2x} \qquad \text{ for } 0 < x
    \]
}

\enquestion{4.21}{
    Determine the cumulative distribution function for the following probability density function:
    \[
        f(x) = \begin{cases}
                    2, & \text{ for } 2.3 < x < 2.9 \\
                    0, & \text{ otherwise.}
        \end{cases}
    \]
    Use the cumulative distribution function to determine the probability that a length exceeds $2.7$ meters.
}

\enquestion{4.29}{
    Suppose that $f(x) = 2x^2$ for $-1 < x < 1$. 
    Determine the mean and variance of $X$.
}

\enquestion{4.37}{
    Integration by parts is required.
    The probability density function of the diameter of a drilled hole in millimeters is $10 e^{-10(x-5)}$ for $x > 5$ mm.
    Although the target diameter is $5$ millimeters, vibrations, tool wear, and other nuisances produce diameters greater than $5$ millimeters.
    \begin{enumerate}[label=(\alph*)]
        \item Determine the mean and variance of the diameter of the holes.
        \item Determine the probability that a diameter exceeds $5.1$ millimeters.
    \end{enumerate}
}

\enquestion{4.41}{
    The thickness of a flange on an aircraft component is uniformly distributed between $0.90$ and $1.05$ millimeters.
    \begin{enumerate}[label=(\alph*)]
        \item Determine the cumulative distribution function of flange thickness.
        \item Determine the proportion of flanges that exceeds 1.02 millimeters.
        \item What thickness is exceeded by \SI{90}{\percent} of the flanges?
        \item Determine the mean and variance of flange thickness.
    \end{enumerate}
    }